\esSezione{Multiplexer 16:1}\label{sec:esercizio1-1}

\begin{traccia}
Progettare, implementare in VHDL e testare mediante simulazione un multiplexer 16:1 a un bit, con sedici ingressi dati e quattro ingressi di selezione, realizzato in modo strutturale a partire da multiplexer 4:1.
\end{traccia}

\progetto

Il multiplexer 16:1 è costruito su due livelli di multiplexer 4:1. Al primo livello ci sono quattro componenti: ognuno riceve un gruppo di quattro ingressi consecutivi e i due bit meno significativi della selezione. Le loro uscite alimentano un quinto multiplexer 4:1, pilotato dai due bit più significativi della selezione, e l'uscita di quest'ultimo è quella del top module. Il top module ha quindi venti ingressi complessivi, sedici dati e quattro di selezione, e un'unica uscita. Lo schema è visibile in figura~\ref{fig:esercizio1-1:schema}.

\begin{figure}[H]
  \centering
  \begin{tikzpicture}[node distance=1cm]
    % primo livello
    \node[blocco, minimum width=2.2cm] (m0) at (0,0)    {\texttt{M0}\\mux\_4\_1};
    \node[blocco, minimum width=2.2cm] (m1) at (0,-2.0) {\texttt{M1}\\mux\_4\_1};
    \node[blocco, minimum width=2.2cm] (m2) at (0,-4.0) {\texttt{M2}\\mux\_4\_1};
    \node[blocco, minimum width=2.2cm] (m3) at (0,-6.0) {\texttt{M3}\\mux\_4\_1};
    % secondo livello
    \node[blocco, minimum width=2.2cm] (m4) at (6.2,-3.0) {\texttt{M4}\\mux\_4\_1};
    % ingressi dati
    \foreach \n in {m0,m1,m2,m3}
      \draw[bus] ($(\n.west)+(-1.9,0)$) -- node[larghezza=4]{/} (\n.west);
    \node[font=\scriptsize\ttfamily, anchor=east] at ($(m0.west)+(-1.9,0)$) {a(0 to 3)};
    \node[font=\scriptsize\ttfamily, anchor=east] at ($(m1.west)+(-1.9,0)$) {a(4 to 7)};
    \node[font=\scriptsize\ttfamily, anchor=east] at ($(m2.west)+(-1.9,0)$) {a(8 to 11)};
    \node[font=\scriptsize\ttfamily, anchor=east] at ($(m3.west)+(-1.9,0)$) {a(12 to 15)};
    % selezione primo livello
    \foreach \n in {m0,m1,m2,m3}
      \draw[bus] ($(\n.south)+(0,-0.6)$) -- node[larghezza=2]{/} (\n.south);
    \foreach \n in {m0,m1,m2,m3}
      \node[font=\scriptsize\ttfamily, anchor=west] at ($(\n.south)+(0.1,-0.4)$) {s(1 downto 0)};
    % uscite primo livello: bus u
    \draw[seg] (m0.east) -| ($(m4.west)+(-0.9,1.0)$) -- ($(m4.west)+(0,1.0)$);
    \draw[seg] (m1.east) -| ($(m4.west)+(-0.7,0.35)$) -- ($(m4.west)+(0,0.35)$);
    \draw[seg] (m2.east) -| ($(m4.west)+(-0.5,-0.35)$) -- ($(m4.west)+(0,-0.35)$);
    \draw[seg] (m3.east) -| ($(m4.west)+(-0.3,-1.0)$) -- ($(m4.west)+(0,-1.0)$);
    \node[font=\scriptsize\ttfamily, anchor=south] at ($(m0.east)+(1.0,0)$) {u(0)};
    \node[font=\scriptsize\ttfamily, anchor=south] at ($(m1.east)+(1.0,0)$) {u(1)};
    \node[font=\scriptsize\ttfamily, anchor=south] at ($(m2.east)+(1.0,0)$) {u(2)};
    \node[font=\scriptsize\ttfamily, anchor=south] at ($(m3.east)+(1.0,0)$) {u(3)};
    % selezione secondo livello e uscita
    \draw[bus] ($(m4.south)+(0,-0.6)$) -- node[larghezza=2]{/} (m4.south);
    \node[font=\scriptsize\ttfamily, anchor=west] at ($(m4.south)+(0.1,-0.4)$) {s(3 downto 2)};
    \draw[seg] (m4.east) -- ++(1.2,0) node[right, font=\scriptsize\ttfamily] {y};
  \end{tikzpicture}
  \caption{MUX 16:1}
  \label{fig:esercizio1-1:schema}
\end{figure}

Nello schema \texttt{u} indica il segnale interno di quattro bit che raccoglie le uscite del primo livello e ne costituisce il vettore di ingresso per \texttt{M4}.

\implementazione

Partiamo dal componente elementare, l'entità \texttt{mux\_4\_1}, che ha come porte \texttt{a}, vettore di quattro ingressi dati indicizzato da 0 a 3, \texttt{s}, selezione a due bit, e \texttt{y}, uscita a un bit. L'architettura \texttt{dataflow} usa un'assegnazione \texttt{with select}: a ogni valore di \texttt{s} corrisponde un ingresso, mentre per i valori non validi l'uscita assume \texttt{'-'}. Il codice è il seguente.

\vhdlfile{esercizio1-1/mux_4_1.vhd}{Implementazione MUX 4:1}{lst:esercizio1-1:mux-4-1}

Definito tale componente, vediamo il top module \texttt{mux\_16\_1}, che ha come porte il vettore \texttt{a} di sedici ingressi (indici da 0 a 15), la selezione \texttt{s} a quattro bit e l'uscita \texttt{y}. L'architettura è \texttt{structural}: il segnale interno \texttt{u}, di quattro bit, raccoglie le uscite dei componenti \texttt{M0}, \texttt{M1}, \texttt{M2} ed \texttt{M3}, istanziati con port mapping sui gruppi \texttt{a(0 to 3)}, \texttt{a(4 to 7)}, \texttt{a(8 to 11)} e \texttt{a(12 to 15)} e con \texttt{s(1 downto 0)} come selezione. Il componente \texttt{M4} riceve \texttt{u} come vettore di ingressi, ha \texttt{s(3 downto 2)} come selezione e produce \texttt{y}. Le istanze usano l'istanziazione diretta dell'entità (\texttt{entity work.mux\_4\_1}), senza dichiarazione di \texttt{component}.

\vhdlfile{esercizio1-1/mux_16_1.vhd}{Implementazione MUX 16:1}{lst:esercizio1-1:mux-16-1}

Rispetto alla soluzione proposta nel modello, l'implementazione differisce nei nomi delle entità (\texttt{mux\_4\_1} e \texttt{mux\_16\_1} al posto di \texttt{mux4\_1} e \texttt{mux16\_1}), nelle istanze, che non passano per una dichiarazione di \texttt{component}, e nel segnale \texttt{u}, privo di valore iniziale.

L'entità \texttt{mux\_4\_1} non verrà ripresentata quando sarà riutilizzata in altri esercizi.

\simulazione

Il testbench \texttt{mux\_16\_1\_tb} ha un'entità priva di porte, poiché non deve essere istanziato. Contiene un'istanza del top module con port mapping sui segnali \texttt{i} (sedici bit, indici da 0 a 15), \texttt{s} (quattro bit) e \texttt{y}. Il process \texttt{stimuli} attende \SI{100}{ns}, assegna a \texttt{i} il valore \texttt{"0010110101001010"} e poi percorre con un ciclo \texttt{for} i valori di \texttt{j} da 0 a 15: a ogni iterazione pone \texttt{s} pari a \texttt{j} (conversione con \texttt{to\_unsigned}), attende \SI{10}{ns} e verifica con un \texttt{assert} che \texttt{y} coincida con \texttt{i(j)}; in caso contrario la simulazione termina con \texttt{severity failure}. Il \texttt{wait} finale sospende il process. Il codice è il seguente.

\vhdlfile{esercizio1-1/mux_16_1_tb.vhd}{Testbench MUX 16:1}{lst:esercizio1-1:mux-16-1-tb}

Il risultato della simulazione è mostrato in figura~\ref{fig:esercizio1-1:simulazione}. L'ingresso \texttt{i} resta costante al valore 0x2D4A ($0010\,1101\,0100\,1010_2$), mentre la selezione \texttt{s} assume in successione i valori da 0x0 a 0xF con un passo di \SI{10}{ns}, a partire da \SI{100}{ns}.

\begin{figure}[H]
  \centering
  \includegraphics[width=0.95\linewidth]{figure/esercizio1-1/simulazione}
  \caption{Simulazione MUX 16:1}
  \label{fig:esercizio1-1:simulazione}
\end{figure}

Poiché l'indice 0 di \texttt{i} corrisponde al bit più a sinistra, con \texttt{s} pari a 0x2 ($0010_2$) l'uscita è \texttt{i(2)}, ossia \texttt{'1'}, e infatti \texttt{y} è alto per tutto l'intervallo in cui la selezione vale 2. Con \texttt{s} pari a 0x3 ($0011_2$) viene selezionato \texttt{i(3)}, uguale a \texttt{'0'}, e \texttt{y} torna basso. Con \texttt{s} pari a 0xF ($1111_2$) viene selezionato \texttt{i(15)}, uguale a \texttt{'0'}, e l'uscita osservata è \texttt{'0'}, coerente con il controllo eseguito dall'\texttt{assert}.
